\documentclass[11pt]{article}
\usepackage[margin=1in]{geometry}
\usepackage{amsmath,amssymb,booktabs,longtable,tabularx,array}
\usepackage{enumitem}
\usepackage[T1]{fontenc}
\usepackage[hidelinks]{hyperref}
\setlist[itemize]{leftmargin=*}
\setlist[enumerate]{leftmargin=*}
\newcommand{\code}[1]{\texttt{\detokenize{#1}}}
\newcommand{\st}{s_t}
\newcommand{\snext}{s_{t+1}}
\newcommand{\obs}{o_t}
\newcommand{\sobs}{\sigma_{\mathrm{obs}}}
\newcommand{\ssafe}{s_{\mathrm{safe}}}
\newcommand{\E}{\mathbb{E}}
\newcommand{\Var}{\mathrm{Var}}
\newcommand{\1}{\mathbf{1}}
\sloppy


\title{State, Action, Reward, and Metric Reference}
\author{Ecology Benchmark Documentation}
\date{July 2026}

\begin{document}
\maketitle

\section{Purpose}

This is the canonical reference for state variables, action semantics, reward
channels, safety thresholds, and evaluator metrics across real, dummy, and
synthetic modes.

\section{State Variables}

\begin{center}
\begin{tabularx}{\linewidth}{llX}
\toprule
Variable & Visibility & Meaning \\
\midrule
$s_t$ & private/evaluator & true continuous abundance; exact zero is terminal \\
$o_t$ & public & log-normal survey observation \\
$r_{\mathrm{base}}$ & private & hidden baseline growth draw for synthetic and shared interface \\
$C$ & private & Allee threshold \\
$\theta$ & private & theta-logistic curvature \\
$z_t$ & private & regime indicator for the switching family \\
$\rho_t$ & public when controls enabled & growth set-point/control state \\
$\kappa_t$ & public when controls enabled & capacity accumulator \\
$K^{\mathrm{eff}}_t$ & public when controls enabled & clipped effective capacity \\
\bottomrule
\end{tabularx}
\end{center}

For set-point ecology cells, the hidden baseline draw is not added to the action
set-point; $\rho$ already carries the growth value used by the transition after
data-derived clipping.

\section{Action Semantics}

Real and dummy ecology cells use set-point growth and cumulative capacity:
\[
  \rho_{t+1} = \Delta r(a_t),\qquad
  \kappa_{t+1} = \kappa_t + \Delta K(a_t).
\]
The code implements this through \code{accumulator_decay_r = 1.0} and
\code{accumulator_decay_K = 0.0}. Capacity is public because it is a manager-owned
consequence of previous actions. The effective growth rate is public in
set-point ecology because the population identity and action table are public.

Synthetic one-step control uses action effects for the next transition only.
Synthetic cumulative-control rows retain public accumulators and support
regression/reference experiments.

\section{Reward}

The real and dummy set-point reward is computed on the true next state and then
logged into the public dataset:
\[
  R_t =
  \alpha\frac{\snext}{\snext+K_{\mathrm{ref}}}
  - c(a_t)
  - P_m\chi_t.
\]
For real cells, $K_{\mathrm{ref}}=K_{\mathrm{base}}$ for the selected population.
The two main reward modes share the benefit and cost terms:
\[
  P_m =
  \begin{cases}
    0, & \code{reward_mode = yield},\\
    \code{collapse_penalty}, & \code{reward_mode = safe}.
  \end{cases}
\]

The penalty indicator is configured separately. Occupancy charges
$\chi_t=\1\{\snext \le \ssafe\}$; crossing charges only a first downward crossing
from above the floor. Real safe-mode defaults to occupancy.

\section{Safety Floors}

The real default safety fraction is depletion-aware:
\begin{itemize}
  \item large/healthy populations use $0.10K_{\mathrm{base}}$;
  \item mid-scale or moderately depleted populations use at least
    $0.20K_{\mathrm{base}}$;
  \item smallest or severely depleted populations use at least
    $0.25K_{\mathrm{base}}$.
\end{itemize}
The MVP threshold is a separate diagnostic, fixed at 50 individuals by default.

\section{Public Reward vs Evaluator Truth}

The public reward is the logged target used by value-fitting methods. It is not a
filter observation. Evaluator truth is kept in the private sidecar and includes
the true state trajectory, safety-entry flags, private structural parameters, and
true reward diagnostics. This split lets all methods train on the same public
information while still allowing paired, truth-aware evaluation.

\section{Metrics}

\begin{center}
\begin{tabularx}{\linewidth}{lX}
\toprule
Metric & Meaning \\
\midrule
operational return & discounted sum of the reward channel used by the method row \\
true return & evaluator truth-channel return; equal to operational return for real set-point reward \\
collapse entry & whether a trajectory entered the below-safety region from above \\
unsafe occupancy & fraction of evaluated steps with $s_t \le \ssafe$ \\
MVP fraction / breach & occupancy and any-hit diagnostics for $s_t \le 50$ \\
persistence & final abundance remains above the safety floor \\
economic cost & accumulated action cost index \\
filter RMSE / log-RMSE & latent-state estimation diagnostics where a filter exists \\
fallback count & deployment fallback usage for guarded methods \\
wall time & filter and planner timing, recorded separately \\
\bottomrule
\end{tabularx}
\end{center}

\section{Mode Differences}

Real and dummy modes share action semantics and reward structure. They differ in
where populations/actions/effects come from: CSVs for real, in-code tables for
dummy. Synthetic modes are retained for stress/reference experiments and may use
older reward conventions in historical configs.

\section*{Verification}

Verified against: \path{src/real_ecology_benchmark/reward.py}, \path{src/real_ecology_benchmark/envs.py}, \path{src/real_ecology_benchmark/evaluator.py}, \path{src/real_ecology_benchmark/gate.py}, \path{src/real_ecology_benchmark/config.py}, \path{src/real_ecology_benchmark/controls.py}, \path{src/real_ecology_benchmark/dataset.py} @ be96c36

\end{document}
